% Independently audited by random_hosts; boundary reduction wording corrected.
\section{A strict six-cut deficit when both shores have large rank}
\label{sec:four-six-large-shores}

The two shores cannot simultaneously realize the critical states of the one-shore relaxation.
The obstruction uses only the matching-augmentation inequalities and the total forest-state bound.

\begin{lemma}[A matching disjoint from two prescribed matchings]
\label{lem:four-six-common-matching}
Any two perfect matchings of six vertices have a common disjoint perfect matching.
Here disjoint means that the two matchings share no edge.
\end{lemma}

\begin{proof}
There are fifteen perfect matchings of six vertices.
A fixed matching has eight disjoint perfect matchings.
To count these, choose the partner of one vertex in four ways, avoiding its prescribed partner.
On the remaining four vertices, exactly one forbidden matching edge remains, leaving two choices.
Thus two sets of disjoint matchings, each of size eight, must intersect.
\end{proof}

Define a limiting six-port state vector to consist of nonnegative path coordinates $a_e$, four-terminal pairing coordinates $u_{S,\alpha}$, and six-terminal pairing coordinates $w_\beta$.
Put
\[
 U_S=\sum_\alpha u_{S,\alpha},\qquad W=\sum_\beta w_\beta,
 \qquad R_A=\sum_{\beta:\,\beta\cap A=\varnothing}w_\beta.
\]
Let $\mathcal P$ be the set of these vectors satisfying
\[
 a_e\le\tfrac12,\qquad U_S\le1,\qquad W\le1,
 \qquad \sum_ea_e+\sum_SU_S+W\le16,
\]
and, for each perfect matching $A=\{e,f,g\}$,
\begin{equation}\label{eq:four-six-limit-augmentation}
 a_e+(U_{e\cup f}-u_{e|f})
     +(U_{e\cup g}-u_{e|g})+R_A\le2.
\end{equation}
The same inequality is required with each of $e,f,g$ as the distinguished edge.
Including the one-added-edge inequalities only reduces $\mathcal P$; the proof below does not need them.

For two state vectors $x,y$, let $\mathcal B(x,y)$ be their normalized crossing expression from the exact six-cut gluing identity.

\begin{proposition}[An absolute gap for two limiting state vectors]
\label{prop:four-six-bilinear-gap}
There is an absolute constant $\delta_6>0$ such that
\[
 \mathcal B(x,y)\le16-\delta_6
 \qquad(x,y\in\mathcal P).
\]
\end{proposition}

\begin{proof}
For one state vector, define
\[
 E(x)=\sum_ea_e^2+\sum_SU_S^2+W^2.
\]
Dropping all incompatible-pairing subtractions and using $2uv\le u^2+v^2$ gives
\[
 \mathcal B(x,y)\le\tfrac12(E(x)+E(y)).
\]
Since $a_e^2\le a_e/2$ and the forest totals lie in $[0,1]$,
\[
 E(x)\le16-\tfrac12\sum_ea_e\le16,
\]
and the same holds for $y$.

Suppose equality $\mathcal B(x,y)=16$ holds.
Both path sums must vanish.
The aggregate bounds then force every one of the fifteen four-terminal totals, and the six-terminal total, to equal one on both shores.
Indeed, sixteen numbers in $[0,1]$ must have squared sum sixteen.
Equality in the dropped four-sector terms also gives
\begin{equation}\label{eq:four-six-disjoint-pairing-supports}
 u^x_{S,\alpha}u^y_{S,\alpha}=0
 \qquad\text{for every }S,\alpha.
\end{equation}

Choose a six-terminal matching $\beta$ with $w^x_\beta>0$, and a matching $\gamma$ with $w^y_\gamma>0$.
They exist because both six-terminal totals equal one.
Lemma~\ref{lem:four-six-common-matching} supplies a matching $A=\{e,f,g\}$ disjoint from both.
Hence $R^x_A>0$ and $R^y_A>0$.

At the proposed equality state, \eqref{eq:four-six-limit-augmentation} becomes
\[
 R_A\le u_{e|f}+u_{e|g},
 \qquad R_A\le u_{e|f}+u_{f|g},
 \qquad R_A\le u_{e|g}+u_{f|g}.
\]
If $R_A>0$, at least two of these three pairing coordinates are positive.
Applying this to both shores requires at least four positive coordinates among three shared slots.
This contradicts~\eqref{eq:four-six-disjoint-pairing-supports}.
Thus equality is impossible.

The set $\mathcal P$ is compact: it is closed, and every coordinate lies in $[0,1]$.
The bilinear function $\mathcal B$ therefore attains its maximum on $\mathcal P\times\mathcal P$.
Its maximum is strictly below sixteen, proving the proposition.
\end{proof}

\begin{theorem}[Two large shores of a cyclic six-cut]
\label{thm:four-six-two-large-shores}
There is an absolute integer $H$ with the following property.
Let $G$ be a connected loopless cubic multigraph with no cyclic cut of size at most five.
Suppose it has a cyclic six-edge cut whose two shores have ranks $h,k\ge H$.
Then, with $r=h+k+5$,
\[
 c(G)<2^{r-1}.
\]
\end{theorem}

\begin{proof}
Both shores are connected.
A cyclic component of either shore must have at least six boundary edges, since its complement contains a cycle.
It therefore uses the entire six-edge cut.
Any further component would have no boundary edge, contradicting the connectedness of $G$.

The cut has distinct endpoints on each shore.
If one vertex has at least two cut incidences, its internal degree is at most one.
Moving it across the cut preserves cycles on both sides and reduces the cut size to at most five.
This would produce a forbidden smaller cyclic cut.
Thus the shores are normalized six-poles.

Their actual normalized state vectors satisfy the matching-augmentation inequalities with errors tending to zero as their ranks tend to infinity.
They also satisfy the aggregate bound
\[
 \sum_ea_e+\sum_SU_S+W\le16+o(1)
\]
from Theorem~\ref{thm:fixed-port-aggregate}.
The individual forest totals are at most one exactly.
Every convergent subsequence of these state vectors therefore has its limit in $\mathcal P$.

Proposition~\ref{prop:four-six-bilinear-gap} and compactness imply that, for sufficiently large $h,k$, the normalized crossing contribution is at most $16-\delta_6/2$.
Otherwise, a sequence violating this bound would have limiting state vectors contradicting the proposition.

Each shore has at most its affine-space size minus one internal cycles.
Their contribution divided by $2^{h+k}$ is at most $2^{-h}+2^{-k}$.
Increase $H$ until this is less than $\delta_6/2$.
The total cycle count is then strictly below $16\,2^{h+k}=2^{r-1}$.
\end{proof}

This theorem leaves bounded-rank six-cut shores unresolved.
The constant gap is established by compactness; no numerical value of $\delta_6$ is asserted here.
